% This chapter is the complete working manuscript for the eleventh seminar talk.
% The final section contains supplementary material outside the default live
% route.

\chapter{Fredholm equations and finite-dimensional reduction}
\label[chapter]{chap:Elliptic_Equations_Kuranishi_Models}

Talk 10 constructed the intrinsic solution stack of a differential equation
and calculated its tangent complex. For
\[
  P(u)=\Delta u+u^2
\]
at its unique solution, that complex is
\[
  [\Cinfty(M)\xrightarrow{\Delta}\Cinfty(M)].
\]
Its terms are infinite-dimensional, but its deformation and obstruction
spaces are both one-dimensional.

This chapter asks how finite-dimensional cohomology becomes an actual
finite-dimensional nonlinear model. Sobolev completion makes Banach calculus
available. Ellipticity makes the linearization Fredholm. Adding a
finite-dimensional obstruction space makes the augmented equation
submersive, and the Banach implicit function theorem then produces a
finite-dimensional manifold carrying an obstruction section.

The obstruction-space construction is the only general reduction in the live
route. Lyapunov--Schmidt coordinates are used only to calculate the running
example and are developed in full in the supplement. The chapter ends with a
precise boundary: the resulting Kuranishi model describes the ordinary smooth
solution germ and its tangent complex, but has not yet been identified with an
open restriction of the intrinsic derived solution stack.

\section{From smooth sections to a Fredholm problem}
\label[section]{sec:Smooth_Sections_To_Fredholm_Problem}

\subsection{Why an analytic completion is needed}
\label[subsection]{sec:Why_Analytic_Completion_Needed}

Let \(E\to M\) and \(F\to M\) be smooth vector bundles over a closed
manifold, and let
\[
  P\colon\mathcal U\subseteq\Gamma(E)\longrightarrow\Gamma(F)
\]
be a nonlinear differential operator of order \(k\). Here \(\mathcal U\) is
an open subset of the smooth section space. Near a solution \(u_0\), its
linearization is a linear differential operator
\[
  D_{u_0}P\colon\Gamma(E)\longrightarrow\Gamma(F)
\]
of order at most \(k\).

The space \(\Gamma(E)\) is naturally a Fr\'echet space. One may measure a
smooth section and all of its derivatives, but no single norm defines this
topology. The ordinary inverse and implicit function theorems for Banach
spaces therefore do not apply directly. There are inverse function theorems
for suitable Fr\'echet spaces, most notably the Nash--Moser theorem, but they
would introduce much more analysis than is needed here.

The standard alternative for an elliptic equation is to work temporarily at
a fixed Sobolev regularity. Choose bundle metrics, a Riemannian metric, and a
connection. For a nonnegative integer \(r\), the Sobolev norm is schematically
\[
  \lVert u\rVert_{H^r}^2
  =\sum_{j=0}^r\lVert\nabla^ju\rVert_{L^2}^2.
\]
The completion of \(\Gamma(E)\) in this norm is a Hilbert space
\(H^r(E)\). Different choices of metrics and connections give equivalent
norms. Moreover,
\[
  \Gamma(E)=\bigcap_{r\geq0}H^r(E).
\]
Thus smooth sections can be recovered by retaining all Sobolev levels.

An operator of order \(k\) loses \(k\) derivatives. Its natural completion
therefore has the form
\begin{equation}
  \label{eq:Sobolev_Extension_Nonlinear_Operator}
  P_l\colon\mathcal U_{l+k}\subseteq H^{l+k}(E)
  \longrightarrow H^l(F),
\end{equation}
not \(H^l(E)\to H^l(F)\). For sufficiently large \(l\), Sobolev embedding
and multiplication theorems imply that \(P_l\) is smooth. In the geometric
setting of jet-defined nonlinear operators, this is part of the functorial
Sobolev-scale construction in Steffens's Lemma~3.4.5
\cite[Lemma~3.4.5]{Steffens_Representability_Elliptic_Moduli}.

\begin{warning}[Completion changes the ambient space]
  \label[warning]{warn:Sobolev_Completion_Changes_Ambient_Space}
  A Sobolev completion introduces nonsmooth sections. It is a device for
  applying Banach calculus, not a claim that the moduli problem should be
  redefined using sections of one arbitrarily chosen regularity. Elliptic
  regularity will be needed to return to smooth solutions.
\end{warning}

\subsection{The linear elliptic theorem}
\label[subsection]{sec:Linear_Elliptic_Fredholm_Theorem}

We recall the functional-analytic property supplied by ellipticity.

\begin{definition}[Fredholm operator]
  \label[definition]{def:Fredholm_Operator}
  A bounded linear operator \(L\colon X\to Y\) between Banach spaces is
  \emph{Fredholm} if
  \[
    \dim\ker L<\infty,
    \qquad
    \operatorname{im}L\subseteq Y\text{ is closed},
    \qquad
    \dim\operatorname{coker}L<\infty.
  \]
  Its \emph{index} is
  \[
    \operatorname{ind}L
    =\dim\ker L-\dim\operatorname{coker}L.
  \]
\end{definition}

\begin{theorem}[Elliptic operators are Fredholm]
  \label[theorem]{thm:Elliptic_Operators_Are_Fredholm}
  Let \(M\) be closed and let
  \[
    L\colon\Gamma(E)\longrightarrow\Gamma(F)
  \]
  be an elliptic linear differential operator of order \(k\). Its Sobolev
  extension
  \[
    L_l\colon H^{l+k}(E)\longrightarrow H^l(F)
  \]
  is Fredholm. Its kernel consists of smooth sections, and its cokernel can be
  represented by a finite-dimensional space of smooth sections of \(F\).
\end{theorem}

This is the principal input from linear elliptic theory. The symbol condition
called ellipticity in the preceding chapter becomes a finite-dimensionality
statement after Sobolev completion. In particular,
\[
  H^0[H^{l+k}(E)\xrightarrow{L_l}H^l(F)]
  =\ker L_l
\]
and
\[
  H^1[H^{l+k}(E)\xrightarrow{L_l}H^l(F)]
  =\operatorname{coker}L_l
\]
are finite-dimensional.

The theorem does not yet give a finite-dimensional nonlinear model. It says
only that the failure of the linearized equation to be invertible is
finite-dimensional. The implicit function theorem will solve all remaining
directions.

\subsection{The nonlinear return to smoothness}
\label[subsection]{sec:Nonlinear_Return_To_Smoothness}

The second quoted input is nonlinear elliptic regularity.

\begin{theorem}[Nonlinear elliptic regularity]
  \label[theorem]{thm:Nonlinear_Elliptic_Regularity}
  Let \(P\) be a nonlinear elliptic differential operator of order \(k\).
  Suppose that \(l\) is sufficiently large, that
  \(u\in H^{l+k}(E)\), and that
  \[
    P_l(u)=f
  \]
  for a smooth section \(f\in\Gamma(F)\). Then \(u\) is smooth.
\end{theorem}

The required lower bound on \(l\) depends on the operator and the dimension
of \(M\). Since we are free to begin at an arbitrarily high Sobolev level, the
precise optimal bound is irrelevant to the reduction. What matters is that a
sufficiently regular weak solution with smooth right-hand side is smooth.
Steffens records exactly this form in Lemma~3.4.6
\cite[Lemma~3.4.6]{Steffens_Representability_Elliptic_Moduli}.

This theorem will be applied not only to \(P_l(u)=0\), but also to an
augmented equation
\[
  P_l(u)+\iota(c)=0,
\]
where \(c\) belongs to a finite-dimensional space and \(\iota(c)\) is a
smooth section. Thus every point of the finite-dimensional manifold produced
below is represented by a smooth section \(u\).

The analytic division of labor is now clear.
\[
  \begin{gathered}
    \text{Sobolev completion}
    \quad\Longrightarrow\quad
    \text{Banach calculus},\\
    \text{ellipticity}
    \quad\Longrightarrow\quad
    \text{Fredholm linearization},\\
    \text{elliptic regularity}
    \quad\Longrightarrow\quad
    \text{smooth solutions}.
  \end{gathered}
\]

\section{The Fredholm--Kuranishi reduction}
\label[section]{sec:Fredholm_Kuranishi_Reduction}

We now isolate the nonlinear construction from its elliptic origin. Let
\(X\) and \(Y\) be Banach spaces, let \(U\subseteq X\) be open, and let
\[
  P\colon U\longrightarrow Y
\]
be smooth. Fix \(x_0\in U\) with \(P(x_0)=0\), and suppose that
\[
  L=D_{x_0}P\colon X\longrightarrow Y
\]
is Fredholm.

\subsection{Splitting the linearized equation}
\label[subsection]{sec:Splitting_Linearized_Equation}

Put \(K=\ker L\). Since \(K\) is finite-dimensional, it admits a closed
complement \(X'\subseteq X\). Since \(\operatorname{im}L\) is closed and has
finite codimension, choose a finite-dimensional subspace \(C\subseteq Y\)
which maps isomorphically to \(\operatorname{coker}L\). We then have
splittings
\begin{equation}
  \label{eq:Fredholm_Splittings}
  X=K\oplus X',
  \qquad
  Y=\operatorname{im}L\oplus C.
\end{equation}
Write
\[
  \iota\colon C\hookrightarrow Y
\]
for the chosen inclusion. The restriction
\begin{equation}
  \label{eq:Fredholm_Complement_Isomorphism}
  L|_{X'}\colon X'\iso\operatorname{im}L
\end{equation}
is an isomorphism of Banach spaces.

The kernel \(K\) consists of infinitesimal solutions. The space \(C\)
records the failure of the linearized equation to be surjective. Neither
splitting in \Cref{eq:Fredholm_Splittings} is canonical. The resulting chart
will accordingly be a presentation rather than an intrinsic object.

\subsection{Adding an obstruction space}
\label[subsection]{sec:Adding_Obstruction_Space}

Consider the augmented map
\begin{equation}
  \label{eq:Augmented_Fredholm_Map}
  \widetilde P\colon U\times C\longrightarrow Y,
  \qquad
  \widetilde P(x,c)=P(x)+\iota(c).
\end{equation}
Its derivative at \((x_0,0)\) is
\[
  D_{(x_0,0)}\widetilde P(v,e)=L(v)+\iota(e).
\]
This map is surjective by the second splitting in
\Cref{eq:Fredholm_Splittings}. Its kernel is
\[
  \ker D_{(x_0,0)}\widetilde P
  =K\times\{0\}.
\]

We use the following form of the Banach implicit function theorem.

\begin{theorem}[Banach submersion theorem]
  \label[theorem]{thm:Banach_Submersion_Theorem}
  Let \(f\colon U\to Y\) be a smooth map between Banach spaces. If
  \(D_xf\) is surjective and has complemented kernel, then, near \(x\), the
  level set \(f^{-1}(f(x))\) is a smooth Banach submanifold with tangent space
  \(\ker D_xf\).
\end{theorem}

In the present application, the kernel is finite-dimensional and hence
complemented. Surjectivity is open for the family of Fredholm operators, so
after shrinking \(U\times C\), we may assume that
\(D\widetilde P\) remains surjective throughout the neighborhood under
consideration.

\begin{theorem}[Fredholm--Kuranishi reduction]
  \label[theorem]{thm:Fredholm_Kuranishi_Reduction}
  Let \(P\colon U\to Y\) and \(x_0\) be as above. After choosing the
  splittings in \Cref{eq:Fredholm_Splittings}, there are a neighborhood
  \(W\subseteq U\times C\) of \((x_0,0)\), a finite-dimensional manifold
  \[
    V=W\cap\widetilde P^{-1}(0),
  \]
  and a smooth map
  \[
    s\colon V\longrightarrow C,
    \qquad
    s(x,c)=c,
  \]
  with the following properties.
  \begin{enumerate}
    \item One has \(\dim V=\dim K\).
    \item The zero locus \(s^{-1}(0)\) is naturally identified with the germ
      of \(P^{-1}(0)\) at \(x_0\).
    \item At every \(v=(x,0)\in s^{-1}(0)\), there is a natural
      quasi-isomorphism
      \begin{equation}
        \label{eq:Kuranishi_Fredholm_Tangent_Comparison}
        [T_vV\xrightarrow{D_vs}C]
        \longrightarrow
        [X\xrightarrow{D_xP}Y].
      \end{equation}
    \item At the base point, the left side of
      \Cref{eq:Kuranishi_Fredholm_Tangent_Comparison} is
      \[
        K\xrightarrow{0}C
      \]
  \end{enumerate}
\end{theorem}

\begin{proof}
  The derivative of \(\widetilde P\) at \((x_0,0)\) is a split surjection
  with kernel \(K\times\{0\}\). By
  \Cref{thm:Banach_Submersion_Theorem}, its zero set is therefore a smooth
  finite-dimensional manifold near \((x_0,0)\), with dimension \(\dim K\).
  This gives \(V\).

  The equation \(s(x,c)=0\) says that \(c=0\). On \(V\), the augmented
  equation also says
  \[
    P(x)+\iota(c)=0.
  \]
  Consequently,
  \[
    s^{-1}(0)
    =\{(x,0)\in W\mid P(x)=0\},
  \]
  which proves the second assertion.

  It remains to compare tangent complexes. At a zero \(v=(x,0)\) of \(s\),
  the tangent space of \(V\) is
  \[
    T_vV
    =\{(\xi,e)\in X\oplus C
      \mid D_xP(\xi)+\iota(e)=0\}.
  \]
  Projection to the two coordinates gives maps
  \[
    \operatorname{pr}_X\colon T_vV\longrightarrow X,
    \qquad
    D_vs=\operatorname{pr}_C\colon T_vV\longrightarrow C.
  \]
  The pair \((\operatorname{pr}_X,-\iota)\) defines the map of complexes in
  \Cref{eq:Kuranishi_Fredholm_Tangent_Comparison}, since
  \[
    D_xP(\operatorname{pr}_X(\xi,e))
    =-\iota(e)
    =-\iota(D_vs(\xi,e)).
  \]

  On degree-zero cohomology, projection to \(X\) gives an isomorphism
  \[
    \ker(D_vs)
    =\{(\xi,0)\mid D_xP(\xi)=0\}
    \iso
    \ker(D_xP).
  \]
  On degree-one cohomology, the induced map is
  \[
    C/\operatorname{im}(D_vs)
    \longrightarrow
    Y/\operatorname{im}(D_xP),
    \qquad
    [e]\longmapsto[-\iota(e)].
  \]
  Surjectivity of
  \[
    D_xP+\iota\colon X\oplus C\longrightarrow Y
  \]
  shows that this map is surjective. If \(\iota(e)=D_xP(\xi)\), then
  \((\xi,-e)\in T_vV\), so \(e\) belongs to
  \(\operatorname{im}(D_vs)\). It is therefore injective as well. The map of
  complexes is a quasi-isomorphism.

  At \((x_0,0)\), one has \(T_{(x_0,0)}V=K\times\{0\}\). Hence
  \(D_{(x_0,0)}s=0\), giving the final assertion.
\end{proof}

The trivial vector bundle \(V\times C\to V\), together with the section
\(s\), is a Kuranishi model in the sense of Chapter~8. Its virtual dimension
is
\begin{equation}
  \label{eq:Fredholm_Index_Virtual_Dimension}
  \dim V-\dim C
  =\dim\ker L-\dim\operatorname{coker}L
  =\operatorname{ind}L.
\end{equation}

\subsection{Returning to smooth elliptic equations}
\label[subsection]{sec:Returning_To_Smooth_Elliptic_Equations}

Apply \Cref{thm:Fredholm_Kuranishi_Reduction} to the Sobolev extension
\(P_l\) in \Cref{eq:Sobolev_Extension_Nonlinear_Operator}. By
\Cref{thm:Elliptic_Operators_Are_Fredholm}, the derivative at a smooth
solution is Fredholm. Choose \(C\) to be represented by smooth sections of
\(F\). Every point \((u,c)\in V\) satisfies
\[
  P_l(u)=-\iota(c),
\]
whose right-hand side is smooth. By
\Cref{thm:Nonlinear_Elliptic_Regularity}, the section \(u\) is smooth.
Consequently, the finite-dimensional manifold \(V\) and its zero locus do not
contain spurious nonsmooth solutions.

We have proved the promised analytic statement.

\begin{corollary}[Local Kuranishi model for an elliptic equation]
  \label[corollary]{cor:Local_Kuranishi_Model_Elliptic_Equation}
  Let \(P\) be a nonlinear elliptic differential operator on a closed
  manifold, and let \(u_0\) be a smooth solution. After choosing a sufficiently
  high Sobolev completion and a finite-dimensional obstruction space, there is
  a finite-dimensional Kuranishi model \((V,V\times C,s)\) such that
  \(s^{-1}(0)\) is the germ of smooth solutions near \(u_0\). At every nearby
  solution, its tangent complex is quasi-isomorphic to the Sobolev completion
  of the linearized differential operator.
\end{corollary}

Linear elliptic regularity identifies the kernel and cokernel of the Sobolev
linearization with those of the operator on smooth sections. Thus the
finite-dimensional tangent complex also recovers the smooth
deformation-obstruction complex from the solution stack in Chapter~10.

\section{The running equation}
\label[section]{sec:Laplacian_Square_Fredholm_Reduction}

We now carry out the reduction completely for the running example. Let \(M\)
be a closed connected Riemannian manifold and let
\[
  P(u)=\Delta u+u^2.
\]
Choose an integer \(l>\dim(M)/2\). Sobolev multiplication makes
\[
  P_l\colon H^{l+2}(M)\longrightarrow H^l(M)
\]
a smooth map. Its derivative at zero is
\[
  D_0P_l=\Delta\colon H^{l+2}(M)\longrightarrow H^l(M).
\]

\subsection{Constants and mean-zero functions}
\label[subsection]{sec:Constants_Mean_Zero_Functions}

Normalize the projection onto constant functions by
\begin{equation}
  \label{eq:Mean_Projection}
  \Pi(f)
  =\frac{1}{\operatorname{vol}(M)}\int_Mf.
\end{equation}
For every Sobolev index \(r\), put
\[
  H^r_0(M)=\ker\Pi.
\]
Then
\begin{equation}
  \label{eq:Constant_Mean_Zero_Splitting}
  H^r(M)=\R\oplus H^r_0(M).
\end{equation}
The kernel of \(\Delta\) is the space of constants, and its image is the
space of mean-zero functions. The restriction
\begin{equation}
  \label{eq:Laplacian_Mean_Zero_Isomorphism}
  \Delta\colon H^{l+2}_0(M)\iso H^l_0(M)
\end{equation}
is an isomorphism. Hence the Fredholm splittings are
\[
  K=\R,
  \qquad
  X'=H^{l+2}_0(M),
  \qquad
  \operatorname{im}\Delta=H^l_0(M),
  \qquad
  C=\R.
\]

Write
\[
  u=a+v,
  \qquad
  a\in\R,
  \quad
  v\in H^{l+2}_0(M).
\]
The mean-zero component of the equation is
\begin{equation}
  \label{eq:Laplacian_Square_Mean_Zero_Equation}
  (1-\Pi)\bigl(\Delta v+(a+v)^2\bigr)=0.
\end{equation}
The derivative of its left side with respect to \(v\) at \((a,v)=(0,0)\) is
the isomorphism in \Cref{eq:Laplacian_Mean_Zero_Isomorphism}. The implicit
function theorem therefore gives a unique local solution
\[
  v=\varphi(a)
\]
of \Cref{eq:Laplacian_Square_Mean_Zero_Equation}.

At this point the example simplifies more than a general
Lyapunov--Schmidt reduction. For every constant \(a\), the choice \(v=0\)
solves the mean-zero equation because
\[
  (1-\Pi)(a^2)=0.
\]
By local uniqueness,
\begin{equation}
  \label{eq:Laplacian_Square_Mean_Zero_Correction_Vanishes}
  \varphi(a)=0
\end{equation}
for all sufficiently small \(a\).

\subsection{The exact obstruction map}
\label[subsection]{sec:Obstruction_Map_Exactly_A_Squared}

The remaining constant component of the equation is
\[
  \kappa(a)
  =\Pi(P_l(a+\varphi(a))).
\]
Using \Cref{eq:Laplacian_Square_Mean_Zero_Correction_Vanishes}, we obtain
\begin{equation}
  \label{eq:Laplacian_Square_Exact_Obstruction_Map}
  \kappa(a)
  =\Pi(a^2)
  =a^2.
\end{equation}
Thus no further coordinate change is needed. The local finite-dimensional
reduction is literally
\[
  \R\longrightarrow\R,
  \qquad
  a\longmapsto a^2.
\]

The same calculation can be displayed using the augmented equation. Include
the constants into \(H^l(M)\) and set
\[
  \widetilde P_l(u,c)=\Delta u+u^2+c.
\]
The mean-zero component again gives \(v=0\). The constant component gives
\[
  a^2+c=0.
\]
Hence the finite-dimensional manifold cut out by the augmented equation is
\[
  V=\{(a,0,-a^2)\mid |a|\text{ is small}\}\cong\R.
\]
Projection to the obstruction coordinate is the section
\[
  s(a)=-a^2.
\]
This is exactly the obstruction-space chart from
\Cref{thm:Fredholm_Kuranishi_Reduction}.

\subsection{The local derived singularity}
\label[subsection]{sec:Laplacian_Square_Local_Derived_Singularity}

The ordinary zero locus of \(a^2\) is one point. Its derived zero locus is
\begin{equation}
  \label{eq:Laplacian_Square_Candidate_Derived_Model}
  \R\times_{\R}^{\mathbf R}\{0\},
  \qquad
  a\longmapsto a^2.
\end{equation}
Its animated \(\Cinfty\)-ring is represented by the Koszul algebra
\[
  \bigl(\Cinfty(\R)\otimes\Lambda(e),\partial e=a^2\bigr),
\]
which was calculated in Chapter~5. Its tangent complex at zero is
\begin{equation}
  \label{eq:Laplacian_Square_Finite_Tangent_Complex}
  [\R\xrightarrow{0}\R].
\end{equation}

The inclusions of constant functions give a map from
\Cref{eq:Laplacian_Square_Finite_Tangent_Complex} to the Sobolev Fredholm
complex
\begin{equation}
  \label{eq:Laplacian_Square_Sobolev_Fredholm_Complex}
  [H^{l+2}(M)\xrightarrow{\Delta}H^l(M)].
\end{equation}
It induces an isomorphism on the kernels and cokernels, hence is a
quasi-isomorphism. Linear elliptic regularity gives a second quasi-isomorphism
to the smooth complex
\[
  [C^\infty(M)\xrightarrow{\Delta}C^\infty(M)]
\]
calculated in Chapter~10.

We have therefore reached the exact finite-dimensional model anticipated at
the beginning of the seminar:
\[
  \begin{gathered}
    \Delta u+u^2=0
    \quad\rightsquigarrow\quad
    a^2=0,\\
    [C^\infty(M)\xrightarrow{\Delta}C^\infty(M)]
    \quad\simeq\quad
    [\R\xrightarrow{0}\R].
  \end{gathered}
\]
The arrow \(\rightsquigarrow\) denotes the classical finite-dimensional
reduction proved in this chapter. Replacing it by an identification of derived
solution stacks is the next task.

\subsection{What the reduction proves, and what remains}
\label[subsection]{sec:Fredholm_Reduction_What_Remains}

The construction in this chapter solves the local analytic problem. It turns
one nonlinear elliptic equation near one solution into a finite-dimensional
section and preserves the entire two-term deformation-obstruction complex.
It also explains the local origin of quasi-smoothness: the candidate derived
model is the zero locus of a section of a finite-rank vector bundle.

Several choices entered the construction:
\begin{enumerate}
  \item A Sobolev index \(l\).
  \item A complement to the kernel of the linearization.
  \item A finite-dimensional subspace representing its cokernel.
  \item A neighborhood on which the augmented operator remains submersive.
\end{enumerate}
Different choices need not give identical Kuranishi charts. This is expected
and was the central lesson of Chapter~8. The intrinsic object should not be
defined to be any one of these presentations.

Steffens instead begins with the solution stack from Chapter~10. The analytic
construction is used to prove that an open restriction of this stack is
represented by the derived zero locus of the obstruction section. Once that
comparison is established, the uniqueness of representing objects supplies
the compatibility among the different local reductions, and the local
representability theorem from Chapter~9 glues them.

The remaining questions are therefore precise.

\begin{enumerate}
  \item Why does the Sobolev reduction compute a pullback in derived smooth
    stacks, rather than only the ordinary solution set?
  \item Why is the result independent of fixing one Sobolev regularity?
  \item How does the construction behave in smooth families of equations?
  \item How do the resulting local representatives glue over an arbitrary
    smooth-stack base?
\end{enumerate}

These questions are the proof architecture of the elliptic representability
theorem.

\section{Summary and supplementary materials}
\label[section]{sec:Elliptic_Kuranishi_Summary_Supplementary_Materials}

\subsection{Takeaways}
\label[subsection]{sec:Fredholm_Reduction_Takeaways}

The live route can be compressed into ten statements.

\begin{enumerate}
  \item Smooth section spaces are Fr\'echet, so the ordinary Banach implicit
    function theorem does not apply directly.
  \item An order-\(k\) differential operator extends at high regularity as
    \(H^{l+k}\to H^l\).
  \item The Sobolev extension of an elliptic linear operator on a closed
    manifold is Fredholm.
  \item Fredholmness says that the failure of invertibility is
    finite-dimensional.
  \item Adding a finite-dimensional representative of the cokernel makes the
    augmented nonlinear equation submersive.
  \item Its zero set is a finite-dimensional manifold carrying an obstruction
    section.
  \item The zero locus of that section is the original solution germ.
  \item The tangent complex of the Kuranishi chart is quasi-isomorphic to the
    Fredholm deformation-obstruction complex.
  \item Elliptic regularity shows that the solutions found after Sobolev
    completion are smooth.
  \item For \(P(u)=\Delta u+u^2\), the obstruction map is exactly
    \(a\mapsto a^2\).
\end{enumerate}

\subsection{Supplement: Lyapunov--Schmidt coordinates}
\label[subsection]{sec:Lyapunov_Schmidt_Derived_Model}

The obstruction-space construction is the form which appears in Steffens's
proof. A second description makes the finite-dimensional obstruction map more
explicit. It also gives the most efficient route through the running example.

\subsubsection{Solving the image equation}
\label[subsubsection]{sec:Solving_Image_Equation}

Retain the splittings
\[
  X=K\oplus X',
  \qquad
  Y=\operatorname{im}L\oplus C
\]
from \Cref{eq:Fredholm_Splittings}. Let
\[
  q\colon Y\longrightarrow\operatorname{im}L,
  \qquad
  \pi_C\colon Y\longrightarrow C
\]
be the two projections. Write a point near \(x_0\) uniquely as
\[
  x=x_0+k+w,
  \qquad
  k\in K,
  \quad
  w\in X'.
\]
The equation \(P(x)=0\) has two components:
\begin{align}
  qP(x_0+k+w)&=0,
    \label{eq:Lyapunov_Schmidt_Image_Equation}\\
  \pi_CP(x_0+k+w)&=0.
    \label{eq:Lyapunov_Schmidt_Obstruction_Equation}
\end{align}

The derivative with respect to \(w\) of the left side of
\Cref{eq:Lyapunov_Schmidt_Image_Equation} at \((k,w)=(0,0)\) is
\[
  qL|_{X'}=L|_{X'}
  \colon X'\iso\operatorname{im}L.
\]
The Banach implicit function theorem therefore gives a unique smooth map
\[
  \varphi\colon K_0\longrightarrow X'
\]
on a neighborhood \(K_0\subseteq K\) of zero such that
\[
  qP(x_0+k+\varphi(k))=0.
\]
The values of \(\varphi\) are the infinite-dimensional directions which have
been solved away.

The remaining equation is finite-dimensional.

\begin{definition}[Lyapunov--Schmidt obstruction map]
  \label[definition]{def:Lyapunov_Schmidt_Obstruction_Map}
  The \emph{Lyapunov--Schmidt obstruction map} is
  \begin{equation}
    \label{eq:Lyapunov_Schmidt_Obstruction_Map}
    \kappa\colon K_0\longrightarrow C,
    \qquad
    \kappa(k)=\pi_CP(x_0+k+\varphi(k)).
  \end{equation}
\end{definition}

By construction, there is an isomorphism of germs of ordinary zero loci
\begin{equation}
  \label{eq:Lyapunov_Schmidt_Zero_Locus_Identification}
  P^{-1}(0)_{x_0}
  \cong
  \kappa^{-1}(0)_0,
  \qquad
  x_0+k+\varphi(k)\longleftrightarrow k.
\end{equation}

The linear term of \(\kappa\) vanishes. Indeed, differentiating the image
equation at zero gives
\[
  L(D_0\varphi(k))=0.
\]
The left side belongs to \(\operatorname{im}L\), while
\(D_0\varphi(k)\in X'\), so
\[
  D_0\varphi=0.
\]
It follows that
\[
  D_0\kappa(k)
  =\pi_CL(k+D_0\varphi(k))
  =0.
\]
Thus all kernel directions are infinitesimal solutions at the base point. The
nonlinear terms of \(\kappa\) determine which of them extend to actual nearby
solutions.

\subsubsection{The two constructions agree}
\label[subsubsection]{sec:Two_Fredholm_Reductions_Agree}

The augmented equation
\[
  P(x)+\iota(c)=0
\]
first implies \(qP(x)=0\). Hence every point of the manifold \(V\) from
\Cref{thm:Fredholm_Kuranishi_Reduction} has the form
\[
  x=x_0+k+\varphi(k).
\]
Its \(C\)-component then says
\[
  \kappa(k)+c=0.
\]
Consequently, the map
\[
  k\longmapsto
  \bigl(x_0+k+\varphi(k),-\kappa(k)\bigr)
\]
is a local parametrization of \(V\), and the obstruction section becomes
\[
  s(k)=-\kappa(k).
\]

Thus the obstruction-space construction and Lyapunov--Schmidt reduction are
not competing methods. They are two descriptions of the same local chart.
The first is coordinate-free once the obstruction space has been chosen. The
second uses kernel and image splittings to write the chart as a map between
finite-dimensional vector spaces.

The sign is harmless. Multiplication by \(-1\) is an automorphism of the
obstruction space, so the derived zero loci of \(\kappa\) and \(-\kappa\) are
isomorphic.

\subsubsection{The tangent complex and virtual dimension}
\label[subsubsection]{sec:Lyapunov_Schmidt_Tangent_Complex}

The Kuranishi chart in Lyapunov--Schmidt coordinates is the section
\[
  \kappa\colon K_0\longrightarrow K_0\times C
\]
of the trivial vector bundle. Its candidate derived zero locus is
\begin{equation}
  \label{eq:Candidate_Derived_Lyapunov_Schmidt_Zero_Locus}
  Z^{\mathrm{der}}(\kappa)
  =K_0\times_C^{\mathbf R}\{0\}.
\end{equation}
At a zero \(k\), its tangent complex is
\[
  [K\xrightarrow{D_k\kappa}C].
\]
By \Cref{thm:Fredholm_Kuranishi_Reduction}, this complex is
quasi-isomorphic to
\[
  [X\xrightarrow{D_{x_k}P}Y],
  \qquad
  x_k=x_0+k+\varphi(k).
\]
At the base point, \(D_0\kappa=0\), so the complex is
\begin{equation}
  \label{eq:Kernel_Cokernel_Tangent_Complex}
  [K\xrightarrow{0}C].
\end{equation}

The quasi-isomorphism is more important than an equality of dimensions. It
identifies the deformation and obstruction spaces and behaves correctly at
nearby solutions where \(D_k\kappa\) need not vanish. Its Euler
characteristic is the Fredholm index by
\Cref{eq:Fredholm_Index_Virtual_Dimension}.

\begin{warning}[What has and has not been proved]
  \label[warning]{warn:Candidate_Derived_Local_Model}
  We have proved that the ordinary solution germ is the zero locus of
  \(\kappa\), and that the tangent complex of its derived zero locus agrees
  with the Fredholm deformation-obstruction complex. These facts strongly
  identify \Cref{eq:Candidate_Derived_Lyapunov_Schmidt_Zero_Locus} as the
  correct local derived model. They do not by themselves prove that this
  derived zero locus represents an open substack of the solution stack from
  Chapter~10.
\end{warning}

The missing assertion concerns the full functor of derived families, not only
its ordinary points and tangent complex. Steffens proves it by showing that
the augmented map is submersive in the derived stack setting. The Sobolev
scale enters that proof as a whole rather than through one fixed completion.


\subsection{Supplement: reduction with parameters}
\label[subsection]{sec:Supplement_Fredholm_Reduction_With_Parameters}

The relative theorem will use a parameterized version of the construction.
Let \(B\) be a finite-dimensional manifold and let
\[
  P\colon B\times U\longrightarrow Y
\]
be smooth. Suppose that \(P(b_0,x_0)=0\) and that the vertical derivative
\[
  L=D_xP(b_0,x_0)\colon X\longrightarrow Y
\]
is Fredholm. Choose \(C\subseteq Y\) representing its cokernel and consider
\[
  \widetilde P(b,x,c)=P(b,x)+\iota(c).
\]
Its derivative in the \((x,c)\)-directions is surjective at
\((b_0,x_0,0)\). Hence the full derivative is surjective, and the zero set
\[
  V=\widetilde P^{-1}(0)
\]
is a finite-dimensional manifold near that point.

The projection \(V\to B\) is a submersion. Indeed, given
\(\dot b\in T_bB\), surjectivity of
\[
  D_xP+\iota\colon X\oplus C\longrightarrow Y
\]
allows us to choose \((\dot x,\dot c)\) satisfying
\[
  D_bP(\dot b)+D_xP(\dot x)+\iota(\dot c)=0.
\]
Then \((\dot b,\dot x,\dot c)\in T_vV\) lifts \(\dot b\).

Projection to \(C\) again gives an obstruction section
\[
  s\colon V\longrightarrow C,
\]
and \(s^{-1}(0)\) is the family of original solutions. At a zero, the
relative tangent complex over \(B\) is quasi-isomorphic to the vertical
Fredholm complex \([X\xrightarrow{D_xP}Y]\).

This finite-dimensional parameterized statement is the local analytic model
behind Steffens's relative theorem. The actual proof must first reduce a
family of section stacks to this form and must then compare the resulting
pullback with the one in derived stacks.

\subsection{Supplement: the Sobolev tower}
\label[subsection]{sec:Supplement_Sobolev_Tower}

A fixed Sobolev completion is enough for the classical reduction, but it is
not an intrinsic replacement for the smooth section stack. Steffens therefore
uses the full tower
\[
  \Gamma(E)\longrightarrow\cdots\longrightarrow
  H^{l+1}(E)\longrightarrow H^l(E)\longrightarrow\cdots.
\]
The transition maps are compact and have dense image, and the smooth section
space is their inverse limit in convenient manifolds. An order-\(k\)
differential operator gives compatible maps
\[
  P_l\colon H^{l+k}(E)\longrightarrow H^l(F).
\]

After adding the same smooth obstruction space at every level, the Banach
implicit function theorem produces finite-dimensional augmented solution
manifolds \(V_l\). Nonlinear elliptic regularity implies that their points are
smooth at sufficiently high regularity. Linear elliptic regularity identifies
their tangent spaces across successive levels. The tower of manifolds
\(V_l\) is therefore eventually locally constant.

This statement contains the analytic heart of Steffens's argument, but it is
not yet the derived conclusion. The Sobolev tower is a limit diagram of smooth
stacks, not a limit diagram of derived \(\Cinfty\)-stacks. The next chapter
will explain how the stacky-submersion criterion bypasses this obstruction.

\subsection{Supplement: a parametrix explanation of Fredholmness}
\label[subsection]{sec:Supplement_Parametrix_Fredholmness}

The proof of \Cref{thm:Elliptic_Operators_Are_Fredholm} belongs to elliptic
analysis, but its mechanism can be summarized. Ellipticity permits the
construction of a pseudodifferential parametrix \(Q\) satisfying
\[
  QL=1-R_1,
  \qquad
  LQ=1-R_2,
\]
where \(R_1\) and \(R_2\) are smoothing operators. On Sobolev spaces over a
closed manifold, smoothing operators are compact. Thus \(L\) is invertible
modulo compact operators, which implies that it is Fredholm.

The same identities improve regularity. If \(Lu=f\) is smooth, then
\[
  u=Qf+R_1u.
\]
Both terms on the right are more regular than the original Sobolev estimate
alone suggests. Iteration gives smoothness. Nonlinear regularity applies this
idea to the linearization along a solution, with estimates controlling the
nonlinear remainder.

This sketch explains why Fredholmness and regularity arise from the same
elliptic mechanism. It is not a proof of the pseudodifferential parametrix
theorem.

\subsection{Supplement: dependence on choices}
\label[subsection]{sec:Supplement_Fredholm_Reduction_Choices}

The finite-dimensional vector spaces \(K\) and \(C\) are determined up to
isomorphism at the base point, but their embeddings and complements are
chosen. Changing a complement changes the solved map \(\varphi\). Enlarging
the obstruction space changes both the dimension of \(V\) and the section
\(s\). Shrinking the neighborhood changes the domain of the chart.

The operations from Chapter~8 already predict the appropriate comparisons.
Shrinking does not change the germ. Adding a redundant obstruction direction
is accompanied by a stabilization of the chart. Changes of complements give
local coordinate changes. These observations explain why different
reductions should present the same derived germ.

They do not replace the representability proof. Constructing comparisons
between every pair of choices, and then constructing all higher
compatibilities, would reproduce the Kuranishi-atlas problem. The intrinsic
solution stack avoids this task: once each reduction is shown to represent an
open restriction of the same stack, their comparisons and higher coherence
follow from uniqueness.

\subsection{Related literature}
\label[subsection]{sec:Fredholm_Reduction_Related_Literature}

The primary source for the chapter is Steffens,
\emph{Representability of Elliptic Moduli Problems in Derived
\(C^\infty\)-Geometry}. Section~1.0.1 explains the passage from Fr\'echet
section spaces to Sobolev completions and separates the analytic problem from
the derived-pullback problem. Lemma~3.4.5 constructs the functorial Sobolev
scale, and Lemma~3.4.6 records the nonlinear regularity theorem
\cite[Section~1.0.1 and Lemmas~3.4.5--3.4.6]{Steffens_Representability_Elliptic_Moduli}.
The proof of Theorem~3.4.3 adds a cokernel representative, applies the Banach
implicit function theorem at every Sobolev level, and then performs the
stack-theoretic comparison which we have postponed
\cite[Proof of Theorem~3.4.3 and
Remark~3.4.7]{Steffens_Representability_Elliptic_Moduli}.

The Banach implicit function theorem may be found in Lang's
\emph{Differential Manifolds}
\cite[Chapter~I]{Lang_Differential_Manifolds}. For the nonlinear elliptic
regularity input, Steffens points to Taylor's
\emph{Pseudodifferential Operators and Nonlinear PDE}, specifically
Theorem~1.2.D
\cite[Theorem~1.2.D]{Taylor_Pseudodifferential_Nonlinear_PDE}.

Pardon begins from the same linear principle: an elliptic complex on a compact
manifold is quasi-isomorphic to a finite-dimensional complex. Section~2 also
explains why this formulation behaves better in families than taking kernels
and cokernels separately
\cite[Section~2]{Pardon_Representability_Elliptic_Fredholm}. In the first step
of the proof sketch of Theorem~5.1, Newton--Picard iteration represents the
regular locus, where the linearization is surjective. The later
finite-dimensional parameter thickening plays the role of adding obstruction
directions
\cite[Proof sketch of Theorem~5.1]{Pardon_Representability_Elliptic_Fredholm}.
